% =============================================================================
%  Capítulo 3 -- Metodología
% =============================================================================
\chapter{Metodología}
\label{cap:metodologia}

\section{Diseño general de la investigación}
\label{sec:diseno-general}

La investigación adopta una metodología de \emph{evaluación sociotécnica} que
combina tres componentes complementarios: (i) un análisis documental, (ii) la
recolección de percepciones sociales mediante encuesta y (iii) un experimento
computacional controlado. El supuesto que articula los tres componentes es
que la evaluación de sistemas algorítmicos desplegados en contextos sociales
concretos no puede limitarse a métricas técnicas abstractas, sino que requiere
contrastar el comportamiento del sistema con el fenómeno social que podría
estar replicando.

\begin{figure}[H]
  \centering
  % Reemplaza este placeholder con un diagrama del flujo metodológico.
  \fbox{\parbox{0.9\textwidth}{\centering\vspace{1cm}
    \textit{[Diagrama del flujo metodológico:}\\
    \textit{revisión documental} $\longrightarrow$ \textit{encuesta de
    percepciones} $\longrightarrow$ \textit{diseño del benchmark} $\longrightarrow$
    \textit{experimentos} $\longrightarrow$ \textit{cálculo del \INS\textit{]}}
    \vspace{1cm}}}
  \caption[Esquema general de la metodología sociotécnica.]{Esquema general
  de la metodología sociotécnica empleada.}
  \label{fig:flujo-metodologico}
\end{figure}

\section{Revisión documental}
\label{sec:revision-documental}

\subsection{Criterios de inclusión y exclusión}

[Bases bibliográficas consultadas; años; idiomas; tipos de documento;
criterios para incluir y excluir referencias.]

\subsection{Estrategia de búsqueda}

[Cadenas de búsqueda, palabras clave en español e inglés, número de resultados
recuperados, filtrado y selección final.]

\subsection{Síntesis temática}

[Descripción del procedimiento para sintetizar la literatura encontrada en
los ejes temáticos del Capítulo~\ref{cap:marco}.]

\section{Recolección de percepciones sobre discriminación tecnológica}
\label{sec:encuesta}

El segundo componente de la metodología consistió en una encuesta dirigida a
personas jóvenes residentes en México, orientada a recoger percepciones y
experiencias sobre el trato diferenciado ---humano y automatizado--- asociado
a la forma de hablar o escribir. La encuesta cumple un doble propósito dentro
del diseño sociotécnico. Por una parte, aporta evidencia situada sobre cómo se
experimenta la discriminación lingüística en el contexto mexicano,
complementando la revisión documental del objetivo~O1. Por otra, sus
respuestas abiertas constituyen el insumo empírico del que se deriva la
tipología de variaciones lingüísticas que estructura el \emph{benchmark}
(Sección~\ref{sec:benchmark-diseno}): en lugar de postular desde el escritorio
qué rasgos del registro escrito son objeto de juicio social, la tipología se
ancla en los rasgos que los propios participantes reportan como sancionados.

\subsection{Población y muestreo}
\label{sec:poblacion-muestreo}

La población objetivo se definió como personas de 18 a 30 años residentes en
México [DECIDIR: confirmar rango de edad con el asesor]. La elección de este
rango respondió a dos criterios. Primero, se trata del grupo con mayor
exposición a procesos automatizados de evaluación de texto ---postulaciones de
empleo y becas en línea, plataformas digitales, sistemas de atención
automatizada---, por lo que sus percepciones refieren a experiencias recientes
y no hipotéticas. Segundo, restringir la población a personas mayores de edad
evita el tratamiento de datos de menores y los requisitos de consentimiento
parental asociados.

El muestreo fue no probabilístico, por conveniencia y con propagación por bola
de nieve: el cuestionario se difundió a través de los contactos del
investigador y se pidió a cada participante compartirlo con otras personas del
rango de edad. Para mitigar el sesgo de red ---la tendencia de este tipo de
muestreo a concentrarse en el círculo inmediato del investigador, en este caso
universitario---, la difusión inicial se dirigió deliberadamente a contactos
de distintos niveles educativos y contextos socioeconómicos
[PENDIENTE: describir brevemente los canales de difusión efectivamente
utilizados].

Esta estrategia de muestreo impone una limitación que conviene dejar asentada
desde ahora: la muestra resultante no es representativa de la población joven
mexicana, por lo que los resultados cuantitativos de la encuesta no admiten
inferencia poblacional. Su función en esta investigación es de carácter
exploratorio e instrumental ---identificar dimensiones relevantes del fenómeno
y alimentar el diseño del \emph{benchmark}---, no estimar prevalencias. Esta
limitación se retoma en la Sección~\ref{sec:limitaciones}.
[PENDIENTE: reportar el tamaño muestral final alcanzado, $n = \ldots$, y el
periodo de levantamiento.]

\subsection{Instrumento}
\label{sec:instrumento}

El instrumento fue un cuestionario en línea, autoadministrado y anónimo, con
una duración estimada de respuesta de 10 a 12 minutos, estructurado en cuatro
secciones: consentimiento informado, datos sociodemográficos, percepciones en
escala tipo Likert y preguntas abiertas. El instrumento completo se reproduce
en el Apéndice~\ref{ap:encuesta}; aquí se describen sus decisiones de diseño.

La sección sociodemográfica (ítems D1--D7) recoge edad, entidad de residencia,
tipo de localidad, nivel educativo propio y situación de actividad. Para
aproximar el nivel socioeconómico de origen se emplearon dos variables
sustitutas (\emph{proxies}): la escolaridad máxima de la madre o tutor
principal (D6) y una autoevaluación retrospectiva de la suficiencia económica
del hogar de crianza (D7). Se optó por estas aproximaciones, de uso común en
encuestas sociales, por ser menos invasivas y menos propensas a la no
respuesta que la pregunta directa por ingreso
[DECIDIR: si el asesor prefiere un instrumento estandarizado, sustituir por la
regla AMAI y ajustar esta justificación].

La sección de percepciones consta de doce ítems (L1--L12) con escala Likert de
cinco puntos, etiquetada en todos sus niveles (de ``totalmente en desacuerdo''
a ``totalmente de acuerdo''), organizados en cuatro bloques: discriminación
lingüística en general (L1--L4), tecnología y trato diferenciado (L5--L6),
escenarios (L7--L9) y experiencia directa (L10--L12). Tres decisiones de
diseño merecen justificación explícita.

Primera, dos ítems (L5 y L7) se redactaron en sentido inverso al resto, de
modo que el acuerdo con ellos expresa la posición contraria a la del bloque.
Su propósito es controlar el sesgo de aquiescencia ---la tendencia a responder
de acuerdo con independencia del contenido--- y detectar patrones de respuesta
automática; ambos se recodifican antes del análisis.

Segunda, el bloque de escenarios sustituye preguntas directas por
mini-viñetas: situaciones concretas (por ejemplo, dos currículums con la misma
experiencia y distinto estilo ante un filtro automatizado) sobre las que el
participante toma postura. Las viñetas reducen la deseabilidad social
---resulta menos costoso emitir un juicio sobre una situación de terceros que
admitir una actitud propia--- y acercan el objeto de la pregunta al escenario
de uso que motiva esta tesis.

Tercera, las cuatro preguntas abiertas (A1--A4) se ordenaron de lo concreto a
lo abstracto: una experiencia vivida de trato diferenciado (A1), las
adaptaciones del propio estilo de escritura ante sistemas automatizados (A2),
el léxico y las formas de escribir que el participante percibe como
penalizadas por una empresa o sistema (A3) y, por último, un mensaje libre a
quienes diseñan estos sistemas (A4). El orden busca anclar primero la memoria
episódica del participante antes de pedirle generalizaciones. El ítem A3
cumple además una función estructural en el diseño: sus respuestas alimentan
directamente la tipología de variaciones del \emph{benchmark}.

Previo al levantamiento, el instrumento se sometió a una prueba piloto con
entre cinco y ocho personas del rango de edad objetivo, con el fin de medir el
tiempo real de respuesta y detectar ítems ambiguos
[PENDIENTE: reportar resultados del piloto y ajustes derivados].

La Tabla~\ref{tab:instrumento-analisis} resume la correspondencia entre las
secciones del instrumento y su uso analítico.

\begin{table}[H]
  \centering
  \caption[Correspondencia entre instrumento y plan de análisis.]{Correspondencia
  entre las secciones del instrumento y su uso en el análisis.}
  \label{tab:instrumento-analisis}
  \begin{tabular}{llp{6.2cm}}
    \toprule
    \textbf{Ítems} & \textbf{Constructo} & \textbf{Uso analítico} \\
    \midrule
    D1--D7   & Perfil sociodemográfico & Caracterización de la muestra \\
    D6--D7   & Nivel socioeconómico (proxy) & Variable de cruce para los bloques Likert \\
    L1--L4   & Discriminación lingüística percibida & Descriptivos y comparación por grupos \\
    L5--L9   & Discriminación tecnológica percibida & Descriptivos y comparación por grupos \\
    L10--L12 & Exposición y confianza & Descriptivos; contexto interpretativo \\
    A1--A2   & Experiencias narradas & Corpus: temas y sentimiento \\
    A3       & Léxico percibido como penalizado & Insumo de la tipología del \emph{benchmark} \\
    A4       & Expectativas normativas & Insumo para la discusión (Cap.~\ref{cap:discusion}) \\
    \bottomrule
  \end{tabular}
\end{table}

\subsection{Análisis cuantitativo}
\label{sec:analisis-cuantitativo}

Las respuestas a los ítems Likert se trataron como variables ordinales. En
consecuencia, el análisis descriptivo reporta distribuciones de frecuencia,
medianas y rangos intercuartílicos por ítem y por bloque, en lugar de medias y
desviaciones estándar, cuya interpretación presupone una escala de intervalo.
Antes del análisis se recodificaron los ítems invertidos (L5 y L7) y se
verificó la consistencia interna de cada bloque mediante el coeficiente alfa
de Cronbach, reportado con carácter descriptivo dada la brevedad de los
bloques [PENDIENTE: reportar valores].

Para explorar si las percepciones varían según el origen socioeconómico, las
respuestas se compararon entre los grupos definidos por las variables proxy
(D6 y D7) mediante pruebas no paramétricas: U de Mann--Whitney para
comparaciones entre dos grupos y H de Kruskal--Wallis para más de dos, con un
nivel de significancia de $\alpha = 0.05$. Dado el carácter no probabilístico
de la muestra, estos contrastes se interpretan como descriptivos de la muestra
obtenida y no como pruebas de hipótesis poblacionales; se reportan junto con
sus tamaños de efecto para facilitar la lectura de la magnitud de las
diferencias [PENDIENTE: completar con los resultados].

\subsection{Análisis del corpus textual}
\label{sec:analisis-corpus}

Las respuestas a las preguntas abiertas se estructuraron en un corpus textual
sobre el que se aplicó la siguiente canalización de procesamiento, implementada
en Python [DECIDIR: fijar versiones de bibliotecas al ejecutar]:

\begin{enumerate}
  \item \textbf{Limpieza y normalización.} Tokenización y lematización con
    \texttt{spaCy} (modelo \texttt{es\_core\_news}), eliminación de palabras
    funcionales y normalización a minúsculas. Una decisión metodológica
    relevante: la normalización se aplicó únicamente sobre la copia de trabajo
    destinada al análisis automático; el texto original de cada respuesta se
    preservó íntegro ---incluidas sus variantes ortográficas---, tanto porque
    esas variantes pueden ser informativas para el objeto de estudio como para
    permitir la cita textual fiel.
  \item \textbf{Análisis de frecuencia.} Frecuencias absolutas y relativas de
    términos y bigramas por pregunta, y ponderación TF--IDF para identificar
    el vocabulario distintivo de cada una.
  \item \textbf{Modelado de temas.} Identificación de temas recurrentes
    mediante asignación latente de Dirichlet (LDA) y, de manera
    complementaria, \texttt{BERTopic} sobre representaciones contextuales en
    español. El número de temas se seleccionó mediante métricas de coherencia
    y revisión manual de interpretabilidad
    [PENDIENTE: reportar configuración y temas resultantes].
  \item \textbf{Análisis de sentimiento.} Clasificación de polaridad con un
    modelo preentrenado para español [DECIDIR: candidato:
    \texttt{pysentimiento}; confirmar al ejecutar], aplicada a las respuestas
    de experiencias (A1--A2) para caracterizar la carga afectiva de las
    narraciones.
\end{enumerate}

El producto final de este análisis es doble: una caracterización temática y
afectiva de las percepciones (reportada en la
Sección~\ref{sec:resultados-encuesta}) y, a partir del ítem A3, un inventario
de rasgos lingüísticos percibidos como penalizados, que constituye el punto de
partida empírico de la tipología de variaciones del \emph{benchmark}
(Sección~\ref{sec:benchmark-diseno}).

\section{Diseño del conjunto de textos controlados}
\label{sec:benchmark-diseno}

\subsection{Principio de control}

El \emph{benchmark} se construye respetando el siguiente principio: dos textos
$t$ y $t'$ del conjunto se consideran un \emph{par controlado} si y sólo si su
contenido informativo es operativamente equivalente y difieren únicamente en
el registro lingüístico. La equivalencia informativa se verifica mediante
revisión humana y, complementariamente, mediante una medida de similitud
semántica sobre representaciones vectoriales.

\subsection{Tipos de variación lingüística considerados}

[Tipología de variaciones a operacionalizar, derivada del análisis de la
encuesta y de la literatura: léxico formal vs.\ coloquial, presencia de
errores ortográficos, uso de anglicismos, longitud oracional, marcadores
discursivos, etc.]

\subsection{Procedimiento de construcción}

[Pasos para construir los pares controlados: redacción inicial, reescritura
con guía estilística, validación por jueces, control de calidad.]

\subsection{Tamaño y composición del \textit{benchmark}}

[Cantidad de textos, número de pares por tipo de variación, distribución de
etiquetas objetivo (en caso de tarea supervisada).]

\section{Evaluación experimental}
\label{sec:experimentos}

\subsection{Modelos evaluados}

[Selección de modelos: al menos un clasificador lineal sobre TF--IDF, al
menos un modelo basado en \textit{transformers} preentrenado en español
(\textit{p.\ ej.,} BETO, RoBERTa en español), y opcionalmente un modelo de
lenguaje grande utilizado como clasificador \textit{zero-shot}.]

\subsection{Tarea de clasificación}

[Definición exacta de la tarea: etiquetas, definiciones operativas, ejemplos.]

\subsection{Protocolo experimental}

\begin{enumerate}
  \item Para cada modelo $f$ y cada par controlado $(t, t')$, se obtienen las
    predicciones $f(t)$ y $f(t')$ y, cuando es posible, las probabilidades
    asociadas $p_f(\cdot \mid t)$ y $p_f(\cdot \mid t')$.
  \item Se calculan métricas estándar de comportamiento agregado por
    \emph{condición de registro}.
  \item Se calcula el \INS según la definición del Capítulo~\ref{cap:ins}.
\end{enumerate}

\subsection{Reproducibilidad}

[Política de versiones de modelos y bibliotecas; semillas aleatorias;
disposición del código y del \emph{benchmark}.]

\section{Consideraciones éticas}
\label{sec:etica}

Las decisiones éticas de esta investigación se concentran en dos frentes: el
tratamiento de las personas participantes en la encuesta y el uso responsable
de los artefactos producidos.

Respecto a las personas participantes, la encuesta se aplicó únicamente a
mayores de edad, con participación voluntaria y posibilidad de abandono en
cualquier momento sin consecuencia alguna. El cuestionario abre con un
consentimiento informado que declara el propósito académico del estudio, la
institución y los responsables, y las condiciones de uso de las respuestas; el
formulario no puede completarse sin aceptarlo expresamente. El instrumento no
solicita nombre, correo, teléfono ni ningún otro dato que permita la
identificación directa de la persona, y las variables sociodemográficas se
recogen en categorías agregadas. Los resultados cuantitativos se reportan
exclusivamente de forma agregada. Las respuestas abiertas pueden citarse
textualmente en este documento, sin atribución; antes de citar cualquier
fragmento se verificó que no contuviera detalles que permitieran la
reidentificación indirecta (nombres de terceros, empresas o lugares
específicos), suprimiéndolos cuando fue necesario
[PENDIENTE: confirmar tras la selección final de citas]. Los datos crudos se
almacenaron en la plataforma del formulario y en una copia local de acceso
restringido al investigador y su asesor, y se conservarán únicamente durante
el desarrollo de la tesis [DECIDIR: confirmar con el asesor si la Facultad
requiere aval de algún comité o registro adicional para encuestas anónimas de
bajo riesgo].

Respecto a los artefactos, el \emph{benchmark} de pares controlados merece una
consideración específica. Por construcción, contiene textos redactados en
registros socialmente estigmatizados; su propósito es auditar sistemas de
clasificación, pero un conjunto así podría, en principio, emplearse en sentido
contrario ---por ejemplo, para entrenar deliberadamente un clasificador que
penalice esos registros. Para reducir ese riesgo, el \emph{benchmark} se
documenta y distribuye con una declaración explícita de uso previsto
(auditoría y evaluación de equidad), siguiendo la práctica de documentación de
conjuntos de datos discutida en la literatura, y los textos que lo integran
son simulados: ningún par proviene de currículums ni de textos reales de
personas. Finalmente, los experimentos del Capítulo~\ref{cap:resultados}
evalúan modelos preentrenados disponibles públicamente y no implican
recolección adicional de datos personales.
